\section{Struttura di gestione del gruppo}

Essendo il gruppo composto da soli due membri, le attività previste in ottica dello \textit{User-Centered Design} sono state suddivise accorpando i ruoli chiave per garantire la copertura di tutte le responsabilità principali del progetto.

È doveroso sottolineare che questa divisione riguarda esclusivamente la \textbf{presa di responsabilità e la direzione} delle singole aree. Durante l'intero ciclo di vita del progetto, infatti, entrambi i membri collaboreranno in maniera attiva, paritaria e costante alla stesura pratica di tutti i documenti e allo sviluppo tecnico di ogni area della piattaforma.

I ruoli di coordinamento sono stati così distribuiti:

\begin{itemize}
    \item \textbf{Manager del Gruppo}: Antonio Walter De Fusco.\\ Ha la responsabilità primaria di coordinare le attività del team, mantenere la visione globale del lavoro e assicurarsi che le varie fasi vengano completate rispettando rigorosamente i tempi prestabiliti (tramite l'uso delle issue e delle milestone).
    
    \item \textbf{Manager di Design}: Antonio Walter De Fusco.\\ Coordina e guida in prima persona l'ideazione e la prototipazione degli aspetti legati al design dell'interfaccia, garantendo la coerenza grafica e l'aderenza metodologica ai task individuati.

    \item \textbf{Manager della Documentazione}: Gabriele di Palma.\\ Coordina la redazione, la strutturazione formale e la supervisione dei vari documenti e report da produrre nel corso dell'iter formativo e progettuale del corso.
    
    \item \textbf{Manager della Valutazione}: Gabriele di Palma.\\ Coordina l'intera catena di valutazione dell'interfaccia. Ricade sotto di lui la responsabilità di gestire il coinvolgimento degli utenti potenziali e di tradurne i feedback in requisiti funzionali ed empirici.
\end{itemize}
